\begin{figure}[H]
\centering
\begin{tikzpicture}
  \draw (0,0) ellipse [x radius=6, y radius=2.5];
  \node at (0,2.1) {$G$};

  \draw plot [smooth cycle, tension=0.9] coordinates {(-4.2,0.5) (-3.2,1.4) (-1.8,0.8) (-2.2,-0.6) (-3.6,-0.7)};
  \node at (-3.1,1) {$G_A$};

  \draw plot [smooth cycle, tension=0.8] coordinates {(2.3,0.9) (4.6,1.2) (4.4,-0.2) (3.2,-0.9) (2.0,-0.2)};
  \node at (4.2,0.8) {$G_B$};

  \node[point] (a1) at (-3.6,0.2) {};
  \node[point] (a2) at (-2.6,0.6) {};
  \node[point] (a3) at (-2.8,-0.4) {};
  \draw[dotted] (a1) -- (a2) -- (a3);

  \node[point] (b1) at (-0.2,0.9) {};
  \node[point] (b2) at (0.8,0.7) {};
  \node[point] (b3) at (1.4,0.2) {};
  \draw[dotted] (b1) -- (b2) -- (b3);
  \draw[dotted] (a2) -- (b1);
  \draw[dotted] (b3) -- (c1);

  \node[point] (c1) at (2.6,0.5) {};
  \node[point] (c2) at (3.8,0.6) {};
  \node[point] (c3) at (3.4,-0.3) {};
  \draw[dotted] (c1) -- (c2) -- (c3);
\end{tikzpicture}
\caption{Example graph where $\TRAV(x,y) = 0$.}
\label{fig:traversal-no-path}
\end{figure}